\documentclass[a4paper,11pt]{article}
\usepackage[utf8]{inputenc}
\usepackage[T1]{fontenc}
\usepackage[ngerman]{babel}
\usepackage{geometry}
\geometry{left=2.5cm, right=2.5cm, top=2.5cm, bottom=2.5cm}
\usepackage{longtable}
\usepackage{booktabs}
\usepackage{array}
\usepackage{xcolor}
\usepackage{colortbl}
\usepackage{hyperref}
\usepackage{fancyhdr}
\usepackage{tabularx}
\usepackage{enumitem}
\usepackage{graphicx}

\definecolor{headerblue}{HTML}{1565C0}
\definecolor{tableheader}{HTML}{E3F2FD}
\definecolor{pass}{HTML}{2E7D32}
\definecolor{warn}{HTML}{F57F17}

\hypersetup{
    colorlinks=true,
    linkcolor=headerblue,
    urlcolor=headerblue
}

\pagestyle{fancy}
\fancyhf{}
\fancyhead[L]{\textbf{INSPECTAIR}}
\fancyhead[R]{Requirements Specification v2.0}
\fancyfoot[C]{\thepage}
\renewcommand{\headrulewidth}{0.4pt}

\setlength{\parindent}{0pt}
\setlength{\parskip}{6pt}

\begin{document}

\begin{center}
{\LARGE\textbf{INSPECTAIR}}\\[0.3cm]
{\Large Requirements Specification}\\[0.3cm]
{\large Luftqualitätsmessgerät mit ESP32-S3}\\[0.5cm]
\begin{tabular}{ll}
\textbf{Projekt:} & InspectAir -- Luftqualitätsmessgerät \\
\textbf{Version:} & 2.0 \\
\textbf{Datum:} & Februar 2026 \\
\textbf{Autoren:} & Maximilian Liebl, Cedric Stroh, Jannis Messer \\
\textbf{Modul:} & Mikrocomputertechnik 1, DHBW Ravensburg \\
\textbf{Dozent:} & Thomas Stingl (Airbus) \\
\end{tabular}
\end{center}

\vspace{0.5cm}
\noindent\textit{Dieses Dokument wurde unter Verwendung von KI-Tools (Claude, Anthropic) zur Überprüfung erstellt.}
\vspace{0.3cm}

\tableofcontents
\newpage

% ============================================================
\section{Einleitung}
% ============================================================

\subsection{Zweck}
Dieses Dokument spezifiziert die Anforderungen an das InspectAir-System, einen Luftqualitätssensor für den Innenraumbereich. Es dient als Grundlage für Design, Implementierung und Test.

\subsection{Scope}
InspectAir ist ein ESP32-S3-basiertes Embedded System zur Überwachung der Raumluftqualität. Das System misst CO\textsubscript{2}, Feinstaub, VOC, Temperatur und Luftfeuchtigkeit und zeigt die Werte auf einem 4-Zoll-Display an.

\subsection{Definitionen}

\begin{longtable}{p{3cm} p{10cm}}
\toprule
\rowcolor{tableheader}
\textbf{Begriff} & \textbf{Definition} \\
\midrule
\endhead
\textbf{CO\textsubscript{2}} & Kohlendioxid -- Indikator für Raumluftqualität \\
\textbf{PM2.5} & Feinstaub mit Partikeldurchmesser $<$2,5\,µm \\
\textbf{VOC} & Volatile Organic Compounds -- flüchtige organische Verbindungen \\
\textbf{UART} & Universal Asynchronous Receiver/Transmitter -- serielle Schnittstelle \\
\textbf{I\textsuperscript{2}C} & Inter-Integrated Circuit -- serieller Datenbus \\
\textbf{SPI} & Serial Peripheral Interface -- synchroner Datenbus \\
\textbf{NDIR} & Non-Dispersive Infrared -- Messprinzip des MH-Z19C CO\textsubscript{2}-Sensors \\
\textbf{FMCW} & Frequency-Modulated Continuous Wave -- Messprinzip des LD2410C Radars \\
\textbf{MOX} & Metal-Oxide -- Sensortyp des SGP40 VOC-Sensors \\
\textbf{Light Sleep} & Energiesparmodus des ESP32 mit reduzierter Stromaufnahme (${\sim}$0,8\,mA), RAM bleibt erhalten, Aufwachzeit ${\sim}$2\,ms \\
\textbf{LVGL} & Light and Versatile Graphics Library -- UI-Framework Version 9.2 \\
\textbf{LovyanGFX} & Display-Treiberbibliothek für ESP32, Version 1.1.16 \\
\bottomrule
\end{longtable}

% ============================================================
\section{Funktionale Anforderungen}
% ============================================================

Die folgenden Anforderungen beschreiben, \textbf{was} das System tun soll.

\begin{longtable}{p{1.8cm} p{2.3cm} p{4cm} p{3.5cm} p{1cm}}
\toprule
\rowcolor{tableheader}
\textbf{ID} & \textbf{Name} & \textbf{Beschreibung} & \textbf{Akzeptanzkriterium} & \textbf{Prio} \\
\midrule
\endhead

\textbf{REQ-F-001} & CO\textsubscript{2}-Messung & Das System muss die CO\textsubscript{2}-Konzentration in ppm messen. & Messbereich 400--5000\,ppm, Genauigkeit $\pm$50\,ppm & Muss \\
\midrule

\textbf{REQ-F-002} & Feinstaub-Messung & Das System muss Feinstaubwerte (PM1.0, PM2.5, PM10) messen. & Messbereich 0--500\,µg/m³, Auflösung 1\,µg/m³ & Muss \\
\midrule

\textbf{REQ-F-003} & VOC-Messung & Das System muss flüchtige organische Verbindungen als Index messen. & VOC-Index 0--500 & Muss \\
\midrule

\textbf{REQ-F-004} & Temperatur-Messung & Das System muss die Umgebungstemperatur messen. & Messbereich 0--50\,°C, Genauigkeit $\pm$0,5\,°C & Muss \\
\midrule

\textbf{REQ-F-005} & Luftfeuchte-Messung & Das System muss die relative Luftfeuchtigkeit messen. & Messbereich 0--100\,\% rH, Genauigkeit $\pm$3\,\% & Muss \\
\midrule

\textbf{REQ-F-006} & Präsenzerkennung & Das System muss Personen im Raum erkennen. & Erkennungsreichweite bis 2\,m, Reaktionszeit $<$500\,ms & Muss \\
\midrule

\textbf{REQ-F-007} & Display-Anzeige & Das System muss alle Messwerte auf einem Display darstellen. & Alle 6 Werte gleichzeitig sichtbar, lesbar aus 1\,m Entfernung & Muss \\
\midrule

\textbf{REQ-F-008} & Farbcodierung & Das System muss Messwerte farblich nach Grenzwerten kategorisieren. & Vierstufige Ampel (Grün/Gelb/Orange/Rot) für CO\textsubscript{2}, PM2.5, VOC & Muss \\
\midrule

\textbf{REQ-F-009} & UI-Wechsel & Der Benutzer muss zwischen mehreren UI-Designs wechseln können. & Button-Druck wechselt zyklisch zwischen 5~Screen-Designs & Muss \\
\midrule

\textbf{REQ-F-010} & Display-Dimming & Das Display muss nach Inaktivität automatisch dimmen. & Dimmen nach 60\,s ohne Präsenz auf 15\,\% Helligkeit & Soll \\
\midrule

\textbf{REQ-F-011} & Display-Aufwecken & Das Display muss bei Präsenzerkennung aufwachen. & Aufwecken innerhalb 500\,ms nach Präsenzerkennung & Soll \\
\midrule

\textbf{REQ-F-012} & Light Sleep & Das System soll bei langer Inaktivität in Stromsparmodus gehen. & Light Sleep nach 5\,Min, Aufwachen durch Radar-Interrupt & Kann \\

\bottomrule
\end{longtable}

% ============================================================
\section{Nicht-Funktionale Anforderungen}
% ============================================================

Die folgenden Anforderungen beschreiben, \textbf{wie gut} das System funktionieren soll.

\begin{longtable}{p{1.8cm} p{2.3cm} p{4cm} p{3.5cm} p{1cm}}
\toprule
\rowcolor{tableheader}
\textbf{ID} & \textbf{Name} & \textbf{Beschreibung} & \textbf{Akzeptanzkriterium} & \textbf{Prio} \\
\midrule
\endhead

\textbf{REQ-NF-001} & Update-Rate & Die Messwerte müssen regelmäßig aktualisiert werden. & Sensor-Update alle 2 Sekunden & Muss \\
\midrule

\textbf{REQ-NF-002} & Reaktionszeit UI & Die Benutzeroberfläche muss flüssig reagieren. & UI-Refresh 10\,Hz, Button-Reaktion $<$200\,ms & Muss \\
\midrule

\textbf{REQ-NF-003} & Dauerbetrieb & Das System muss für 24/7-Dauerbetrieb geeignet sein. & Keine Abstürze oder Memory Leaks über 24\,h & Muss \\
\midrule

\textbf{REQ-NF-004} & Stromversorgung & Das System muss über USB betrieben werden können. & USB 5\,V, max. 1\,A Stromaufnahme & Muss \\
\midrule

\textbf{REQ-NF-005} & Gehäuse & Das System soll in einem 3D-gedruckten PETG-Gehäuse untergebracht werden. & Belüftungsöffnungen für Sensoren, Display-Ausschnitt, USB-Zuführung & Soll \\
\midrule

\textbf{REQ-NF-006} & Wartbarkeit & Der Code soll modular und dokumentiert sein. & Mind. 80\,\% der Funktionen mit Doxygen dokumentiert, max. 3 Module-Dependencies & Soll \\

\bottomrule
\end{longtable}

% ============================================================
\section{Hardware-Anforderungen}
% ============================================================

\begin{longtable}{p{1.8cm} p{2.5cm} p{4cm} p{4.5cm}}
\toprule
\rowcolor{tableheader}
\textbf{ID} & \textbf{Komponente} & \textbf{Beschreibung} & \textbf{Akzeptanzkriterium} \\
\midrule
\endhead

\textbf{REQ-HW-001} & Mikrocontroller & ESP32-S3-WROOM N8R8 DevKit (Waveshare) & Dual-Core, WiFi/BLE, 8\,MB Flash QIO\_OPI, 8\,MB PSRAM \\
\midrule
\textbf{REQ-HW-002} & Display & 4-Zoll TFT-Display mit ST7796S Controller & 480$\times$320 Pixel, SPI-Interface \\
\midrule
\textbf{REQ-HW-003} & CO\textsubscript{2}-Sensor & MH-Z19C NDIR CO\textsubscript{2}-Sensor & SoftwareSerial 9600\,Baud, 5\,V Versorgung \\
\midrule
\textbf{REQ-HW-004} & Feinstaubsensor & PMS5003 Laser-Partikelsensor & UART1 9600\,Baud, 5\,V Versorgung \\
\midrule
\textbf{REQ-HW-005} & VOC-Sensor & SGP40 MOX VOC-Sensor & I\textsuperscript{2}C 0x59, 3,3\,V \\
\midrule
\textbf{REQ-HW-006} & Temp/Feuchte-Sensor & AHT20 Temperatur- und Feuchtigkeitssensor & I\textsuperscript{2}C 0x38, 3,3\,V \\
\midrule
\textbf{REQ-HW-007} & Radar-Sensor & LD2410C 24\,GHz mmWave Radar & UART2 256000\,Baud, Level Shifter erforderlich! \\
\midrule
\textbf{REQ-HW-008} & Level Shifter & Bidirektionaler Level Shifter (TXS0108E) & 5\,V$\leftrightarrow$3,3\,V Wandlung für UART TX \\
\midrule
\textbf{REQ-HW-009} & Stützkondensatoren & Elkos zur Pufferung von Stromspitzen & 220\,µF 25\,V an MH-Z19C und LD2410C VCC \\
\midrule
\textbf{REQ-HW-010} & MOSFET & BS170 N-Kanal MOSFET für Backlight-PWM & Logic-Level, Gate über 1\,k$\Omega$ an GPIO3 \\
\midrule
\textbf{REQ-HW-011} & Button & Taster für UI-Wechsel & GPIO1 mit internem Pull-up \\

\bottomrule
\end{longtable}

% ============================================================
\section{Grenzwert-Definition}
% ============================================================

Die Farbcodierung der Messwerte basiert auf folgenden Grenzwerten (implementiert in \texttt{colors.h}):

\begin{longtable}{p{2.5cm} p{2.5cm} p{2.5cm} p{2.5cm} p{2.5cm}}
\toprule
\rowcolor{tableheader}
\textbf{Parameter} & \textbf{Gut (Grün)} & \textbf{Mäßig (Gelb)} & \textbf{Schlecht (Orange)} & \textbf{Gefährlich (Rot)} \\
\midrule
\endhead

\textbf{CO\textsubscript{2}} & $<$ 800\,ppm & 800--1000\,ppm & 1000--1500\,ppm & $>$ 1500\,ppm \\
\midrule
\textbf{PM2.5} & $\leq$ 5\,µg/m³ & 5--15\,µg/m³ & 15--25\,µg/m³ & $>$ 25\,µg/m³ \\
\midrule
\textbf{VOC Index} & $\leq$ 100 & 100--200 & 200--300 & $>$ 300 \\
\midrule
\textbf{Temperatur} & 18--26\,°C & 15--18 / 26--30\,°C & \multicolumn{2}{l}{$<$ 15 / $>$ 30\,°C} \\
\midrule
\textbf{Luftfeuchte} & 30--60\,\% & 20--30 / 60--70\,\% & \multicolumn{2}{l}{$<$ 20 / $>$ 70\,\%} \\

\bottomrule
\end{longtable}

\textit{Hinweis:} Die Grenzwerte für PM2.5 orientieren sich an den WHO-Richtlinien 2021 und wurden gegenüber Version~1.0 verschärft. Die Farbcodierung nutzt vier Stufen (Grün/Gelb/Orange/Rot) statt der ursprünglichen drei Stufen.

% ============================================================
\section{Requirements Traceability}
% ============================================================

Jede Anforderung wird durch entsprechende Testfälle verifiziert:

\begin{longtable}{p{2.2cm} p{4cm} p{6cm}}
\toprule
\rowcolor{tableheader}
\textbf{Requirement} & \textbf{Beschreibung} & \textbf{Testfall} \\
\midrule
\endhead

\textbf{REQ-F-001} & CO\textsubscript{2}-Messung & TC-04 (Einzeltest), TC-09 (Paralleltest) \\
\textbf{REQ-F-002} & Feinstaub-Messung & TC-03 (Einzeltest), TC-09 (Paralleltest) \\
\textbf{REQ-F-003} & VOC-Messung & TC-02 (Einzeltest), TC-09 (Paralleltest) \\
\textbf{REQ-F-004} & Temperatur-Messung & TC-02 (Einzeltest), TC-09 (Paralleltest) \\
\textbf{REQ-F-005} & Luftfeuchte-Messung & TC-02 (Einzeltest), TC-09 (Paralleltest) \\
\textbf{REQ-F-006} & Präsenzerkennung & TC-05 (Einzeltest), TC-15 (Wake-Test) \\
\textbf{REQ-F-007} & Display-Anzeige & TC-01 (Display), TC-10, TC-11 \\
\textbf{REQ-F-008} & Farbcodierung & TC-10, TC-11, TC-13 (UI Consistency) \\
\textbf{REQ-F-009} & UI-Wechsel & TC-12 (Button-Wechsel) \\
\textbf{REQ-F-010} & Display-Dimming & TC-07 (MOSFET), TC-14 (Dimming) \\
\textbf{REQ-F-011} & Display-Aufwecken & TC-15 (Wake-Test) \\
\textbf{REQ-F-012} & Light Sleep & TC-16 (Light Sleep) \\
\textbf{REQ-HW-009} & Stützkondensatoren & TC-08 (Stabilitätstest) \\

\bottomrule
\end{longtable}

% ============================================================
\section*{Änderungshistorie}
% ============================================================

\begin{longtable}{l l l p{7cm}}
\toprule
\rowcolor{tableheader}
\textbf{Version} & \textbf{Datum} & \textbf{Autor} & \textbf{Änderung} \\
\midrule
1.0 & Januar 2026 & Team & Erstversion \\
2.0 & Februar 2026 & Team & Gehäuse: 3D-Druck PETG statt Holz (REQ-NF-005). Pin-Belegung aktualisiert gemäß Code. Grenzwerte aktualisiert auf WHO 2021 mit 4-Stufen-Ampel gemäß \texttt{colors.h}. UI-Wechsel erweitert auf 5~Screens. Deep Sleep durch Light Sleep ersetzt. \\
\bottomrule
\end{longtable}

\end{document}
